% !TEX program = xelatex
% ALW 实验进展汇报（16:9）。编译：见 README.md（WSL 上用 tools/render.py；Overleaf 选 XeLaTeX）。
% 数字全部来自 facts.tex（tools/gen_facts.py 生成）；今晚的结果只在 tonight_slots.tex 里填。
\documentclass[aspectratio=169,10pt,t]{ctexbeamer}
% theme.tex -- look, components and the fact / slot interfaces of the ALW progress deck.
% Packages are limited to the offline-cached set: ctexbeamer, tikz, booktabs, array, multirow.
\usepackage{booktabs}
\usepackage{array}
\usepackage{multirow}
\usepackage{tikz}
\usetikzlibrary{positioning,calc,fit,arrows.meta}

% ---------- fonts: Microsoft YaHei + Consolas on WSL; Fandol / Latin Modern elsewhere ----------
\IfFileExists{/mnt/c/Windows/Fonts/msyh.ttc}{%
  \setCJKmainfont{msyh.ttc}[Path=/mnt/c/Windows/Fonts/,BoldFont=msyhbd.ttc,ItalicFont=msyh.ttc]
  \setCJKsansfont{msyh.ttc}[Path=/mnt/c/Windows/Fonts/,BoldFont=msyhbd.ttc,ItalicFont=msyh.ttc]
  \setCJKmonofont{msyh.ttc}[Path=/mnt/c/Windows/Fonts/,BoldFont=msyhbd.ttc]
  \setsansfont{msyh.ttc}[Path=/mnt/c/Windows/Fonts/,BoldFont=msyhbd.ttc,ItalicFont=msyh.ttc]
}{}
\IfFileExists{/mnt/c/Windows/Fonts/consola.ttf}{%
  \setmonofont{consola.ttf}[Path=/mnt/c/Windows/Fonts/,BoldFont=consolab.ttf,Scale=0.92]
}{}

% ---------- colours (meaning is fixed across the deck) ----------
\definecolor{accent}{RGB}{31,78,121}     % what we added
\definecolor{ink}{RGB}{34,40,49}
\definecolor{muted}{RGB}{108,114,124}    % untouched benchmark / Claude Code
\definecolor{hair}{RGB}{212,217,224}
\definecolor{paper}{RGB}{245,247,250}
\definecolor{ok}{RGB}{34,124,82}         % pass
\definecolor{bad}{RGB}{188,56,50}        % fail
\definecolor{slot}{RGB}{186,118,10}      % tonight's slot
\definecolor{lock}{RGB}{70,76,86}        % sealed val / test, blind aggregates

% ---------- beamer ----------
\setbeamertemplate{navigation symbols}{}
\setbeamersize{text margin left=8mm,text margin right=8mm}
\setbeamercolor{normal text}{fg=ink}
\setbeamercolor{structure}{fg=accent}
\setbeamercolor{frametitle}{fg=accent}
\setbeamerfont{frametitle}{size=\large,series=\bfseries}
\setbeamertemplate{frametitle}{%
  \vskip3.6mm%
  {\color{accent}\rule[-1.1mm]{1.1mm}{5mm}}\hspace{2.4mm}%
  \usebeamerfont{frametitle}\usebeamercolor[fg]{frametitle}\insertframetitle\par\vskip-0.6mm}
\setbeamertemplate{footline}{%
  \hbox to\paperwidth{\hspace{8mm}{\tiny\color{muted}\insertshorttitle}\hfill
    {\scriptsize\color{muted}\insertframenumber}\hspace{8mm}}\vskip2.2mm}
\setbeamertemplate{itemize items}{\color{accent}\rule[0.25ex]{0.9ex}{0.9ex}}
\setlength{\leftmargini}{4.2mm}
\setbeamerfont{itemize/enumerate body}{size=\small}

\makeatletter
% ---------- small text components ----------
\newcolumntype{L}[1]{>{\raggedright\arraybackslash}p{#1}}
\newcolumntype{C}[1]{>{\centering\arraybackslash}p{#1}}
\newcolumntype{R}[1]{>{\raggedleft\arraybackslash}p{#1}}
\newcommand{\hl}[1]{\textcolor{accent}{\textbf{#1}}}
\newcommand{\code}[1]{\texttt{#1}}
\newcommand{\Ystar}{$Y^\star$}
\newcommand{\Cstar}{$C^\star$}
\newcommand{\yes}{\textcolor{ok}{是}}
\newcommand{\no}{\textcolor{muted}{否}}
\newcommand{\pass}{\textcolor{ok}{\textbf{通过}}}
\newcommand{\fail}{\textcolor{bad}{\textbf{未通过}}}
\newcommand{\lead}[1]{{\small #1\par}\vspace{1.6mm}}
% evidence tags
\newcommand{\evtag}[2]{\textcolor{#1}{\textbf{[#2]}}}
\newcommand{\tagF}{\evtag{ok}{事实}}
\newcommand{\tagJ}{\evtag{accent}{判断}}
\newcommand{\tagU}{\evtag{slot}{未核验}}
\newcommand{\tagM}{\evtag{muted}{机制}}
\newcommand{\tagE}{\evtag{muted}{估算}}
\newcommand{\tagS}{\evtag{slot}{待填}}
% source line at the bottom of a frame
\newcommand{\src}[1]{\par\vfill{\scriptsize\color{muted}#1\par}\vspace{0.4mm}}
\newcommand{\pill}[2]{{\setlength{\fboxsep}{1.6pt}\colorbox{#1}{\textcolor{white}{\scriptsize\bfseries\,#2\,}}}}
\newcommand{\chip}[2]{{\setlength{\fboxsep}{1.6pt}\colorbox{#1!13}{\textcolor{#1}{\scriptsize\,#2\,}}}}

% ---------- facts: \setfact{key}{value} (facts.tex, generated) and \fv{key} ----------
\newcommand{\setfact}[2]{\expandafter\gdef\csname fact@#1\endcsname{#2}}
\newcommand{\fv}[1]{\ifcsname fact@#1\endcsname\csname fact@#1\endcsname\else
  \PackageError{deck}{Unknown fact '#1' (regenerate facts.tex)}{}\fi}

% ---------- slots: \slot{key}{value} (tonight_slots.tex) ----------
\newcommand{\slot}[2]{\expandafter\gdef\csname slot@#1\endcsname{#2}}
\newcommand{\slotrow}[5]{\csgappto{slotrows@#1}{#2 & #3 & \deltafmt{#4} & #5 \\}}
\newcommand{\ifslot}[3]{\ifcsvoid{slot@#1}{#3}{#2}}                 % {key}{filled}{empty}
\newcommand{\ifslotis}[4]{\ifcsstring{slot@#1}{#2}{#3}{#4}}          % {key}{string}{yes}{no}
\newcommand{\slotraw}[1]{\ifcsvoid{slot@#1}{}{\csname slot@#1\endcsname}}
\newcommand{\slotchip}[1][明早填入]{\chip{slot}{#1}}
\newcommand{\slotval}[2][明早填入]{\ifslot{#2}{\csname slot@#2\endcsname}{\slotchip[#1]}}
% signed number coloured by sign; empty -> chip
\newcommand{\deltafmt}[1]{\edef\deck@tmp{#1}\ifx\deck@tmp\@empty\slotchip\else
  \pgfmathtruncatemacro\deck@sgn{\deck@tmp>0.000001?1:(\deck@tmp<-0.000001?2:0)}%
  \ifcase\deck@sgn\relax\textcolor{muted}{\deck@tmp}\or\textcolor{ok}{\textbf{\deck@tmp}}\or
  \textcolor{bad}{\textbf{\deck@tmp}}\fi\fi}

% ---------- diff card lines (slots/r<g>_patch.tex is generated by tools/fill_slots.py) ----------
\newcommand{\deck@dline}[3]{{\setlength{\fboxsep}{0pt}\colorbox{#1}{\makebox[\linewidth][l]{%
  \color{#2}\ttfamily\rule[-0.28em]{0pt}{1.12em}\,#3}}}\par}
\newcommand{\dadd}[1]{\deck@dline{ok!13}{ok!75!black}{#1}}
\newcommand{\ddel}[1]{\deck@dline{bad!11}{bad!85!black}{#1}}
\newcommand{\dctx}[1]{\deck@dline{white}{ink!85}{#1}}
\newcommand{\dhunk}[1]{\deck@dline{accent!8}{accent!85}{#1}}
\newcommand{\dfile}[1]{\deck@dline{accent!20}{accent}{\bfseries #1}}
\newcommand{\dmore}[1]{\deck@dline{white}{muted}{\rmfamily …… #1}}
\newcommand{\nl}[1]{{\ttfamily #1}\par}

% ---------- TikZ components ----------
\tikzset{
  every picture/.style={x=1mm,y=1mm,line cap=round,line join=round},
  box/.style={draw=hair,fill=paper,rounded corners=1.1mm,inner sep=1.6mm,align=left,
              font=\footnotesize,line width=0.5pt},
  ours/.style={box,draw=accent!70,fill=accent!7},
  base/.style={box,draw=muted!45,fill=muted!7},
  good/.style={box,draw=ok!70,fill=ok!8},
  fail/.style={box,draw=bad!70,fill=bad!7},
  slotbox/.style={box,draw=slot,dashed,fill=slot!6},
  locked/.style={box,draw=lock!70,fill=lock!11,text=lock},
  arr/.style={-{Stealth[length=1.9mm,width=1.5mm]},line width=0.55pt,draw=muted},
  lab/.style={font=\scriptsize,text=muted,inner sep=0.6mm},
  num/.style={circle,fill=accent,text=white,font=\scriptsize\bfseries,inner sep=0pt,
              minimum size=3.9mm},
  tile/.style={box,align=center,inner sep=1.8mm},
}
\newcommand{\bignum}[2][accent]{{\fontsize{19}{21}\selectfont\bfseries\color{#1}#2}}
\newcommand{\midnum}[2][accent]{{\fontsize{13}{15}\selectfont\bfseries\color{#1}#2}}

% visible Phase 1c grid cells (rows from facts.tex)
\newcommand{\tasklabel}[3]{\node[anchor=west,font=\scriptsize,inner sep=0pt] at (0,-#1*4.15)
  {\makebox[6.2mm][l]{\textbf{#2}}#3};}
\newcommand{\runcell}[4]{%
  \def\deck@c{muted}\def\deck@s{#4}%
  \def\deck@k{nogds}\ifx\deck@s\deck@k\def\deck@c{bad}\fi
  \def\deck@k{drc}\ifx\deck@s\deck@k\def\deck@c{accent}\fi
  \def\deck@k{lvs}\ifx\deck@s\deck@k\def\deck@c{ok}\fi
  \node[draw=\deck@c!60,fill=\deck@c!14,text=\deck@c!85!black,font=\scriptsize\bfseries,
        minimum width=8mm,minimum height=3.5mm,inner sep=0pt,rounded corners=0.5mm]
        at (37.5+#2*9,-#1*4.15) {#3};}
\makeatother

% AUTO-GENERATED by tools/gen_facts.py -- do not edit by hand.
% Sources: pack = verified_visible_facts.json; config = configs/alw/config.server.json;
%          log = docs/alw-deployment-log.md (regex-parsed, cross-checked).
% Blind clusters appear only as aggregate counts.

\setfact{split.total}{40}% pack
\setfact{split.dev}{12}% pack
\setfact{split.devVisible}{8}% pack
\setfact{split.blindClusters}{3}% pack
\setfact{split.blindTasks}{4}% pack
\setfact{split.val}{14}% pack
\setfact{split.test}{14}% pack
\setfact{wm.tMax}{10800}% config
\setfact{wm.tWarn}{8400}% config
\setfact{wm.tSoft}{9900}% config
\setfact{wm.repSec}{900}% config
\setfact{wm.repTurns}{80}% config
\setfact{wm.minRemain}{180}% config
\setfact{wm.verifierSec}{3600}% config
\setfact{wm.tMaxMin}{180}% config
\setfact{wm.tWarnMin}{140}% config
\setfact{wm.tSoftMin}{165}% config
\setfact{wm.repMin}{15}% config
\setfact{wm.verifierMin}{60}% config
\setfact{res.cpus}{6}% config
\setfact{res.memGiB}{32}% config
\setfact{agent.slots}{8}% config
\setfact{agent.version}{2.1.302}% config
\setfact{agent.model}{deepseek-v4-1-flash}% log
\setfact{pc.n}{24}% pack
\setfact{pc.nInfra}{0}% pack
\setfact{pc.meanY}{0.1640}% pack
\setfact{pc.sigmaSeed}{0.0223}% pack
\setfact{pc.seedMeanZero}{0.1797}% pack
\setfact{pc.seedMeanOne}{0.1482}% pack
\setfact{pc.taxMirror}{19}% pack
\setfact{pc.taxElec}{3}% pack
\setfact{pc.taxGds}{1}% pack
\setfact{pc.taxCap}{0}% pack
\setfact{pc.nNotPassed}{23}% pack
\setfact{conc.gds.agree}{23}% pack
\setfact{conc.gds.excluded}{1}% pack
\setfact{conc.gds.falseReject}{0}% pack
\setfact{conc.gds.falseCertify}{0}% pack
\setfact{conc.lvs.agree}{23}% pack
\setfact{conc.lvs.excluded}{1}% pack
\setfact{conc.lvs.falseReject}{0}% pack
\setfact{conc.lvs.falseCertify}{0}% pack
\setfact{conc.shape.agree}{21}% pack
\setfact{conc.shape.excluded}{1}% pack
\setfact{conc.shape.falseReject}{0}% pack
\setfact{conc.shape.falseCertify}{2}% pack
\setfact{conc.fill.agree}{5}% pack
\setfact{conc.fill.excluded}{19}% pack
\setfact{conc.fill.falseReject}{0}% pack
\setfact{conc.fill.falseCertify}{0}% pack
\setfact{conc.morph.agree}{5}% pack
\setfact{conc.morph.excluded}{19}% pack
\setfact{conc.morph.falseReject}{0}% pack
\setfact{conc.morph.falseCertify}{0}% pack
\setfact{conc.drc.agree}{23}% pack
\setfact{conc.drc.excluded}{1}% pack
\setfact{conc.drc.falseReject}{0}% pack
\setfact{conc.drc.falseCertify}{0}% pack
\setfact{vis.n}{16}% pack
\setfact{vis.meanY}{0.0938}% pack
\setfact{vis.gds}{15}% pack
\setfact{vis.drc}{9}% pack
\setfact{vis.lvs}{0}% pack
\setfact{vis.promoted}{4}% pack
\setfact{vis.gate}{16}% pack
\setfact{vis.promotedStillLow}{4}% pack
\setfact{vis.wallMin}{5599}% pack
\setfact{vis.wallMax}{10689}% pack
\setfact{blind.nRuns}{8}% pack
\setfact{blind.meanY}{0.304}% pack
\setfact{fun.gds}{23}% log
\setfact{fun.drc}{15}% log
\setfact{fun.lvs}{5}% log
\setfact{fun.pex}{5}% log
\setfact{fun.postsim}{1}% log
\setfact{pc.tsoft}{14}% log
\setfact{pc.gate}{20}% log
\setfact{pc.promoted}{5}% log
\setfact{pc.promotedPct}{25.0}% log
\setfact{pc.disconnects}{13}% log
\setfact{pc.disconnectRuns}{11}% log
\setfact{pc.requests}{9\,732}% log
\setfact{pc.requestsPerRun}{406}% log
\setfact{pc.inputTokM}{755.8}% log
\setfact{pc.outputTokM}{34.4}% log
\setfact{pc.inputTokPerRunM}{31.5}% log
\setfact{pc.outputTokPerRunM}{1.43}% log
\setfact{dist.zero}{1}% log
\setfact{dist.tenth}{18}% log
\setfact{dist.fifth}{1}% log
\setfact{dist.above}{4}% log
\setfact{blind.cSum}{1}% log
\setfact{blind.gds}{8}% log
\setfact{blind.drc}{6}% log
\setfact{blind.lvs}{5}% log
\setfact{blind.pex}{5}% log
\setfact{blind.postsim}{1}% log
\setfact{blind.promoted}{1}% log
\setfact{pbTwo.sec}{200.8}% pack
\setfact{pbTwo.budget}{900.0}% pack
\setfact{pbTwo.pct}{22.3}% pack
\setfact{pbTwo.frac}{0.2232}% pack
\setfact{pbTwo.budgetFrac}{1.0000}% pack
\setfact{pbTwo.rc}{1}% pack
\setfact{pbTwo.turns}{15}% pack
\setfact{pbTwentyone.sec}{167.8}% pack
\setfact{pbTwentyone.budget}{779.9}% pack
\setfact{pbTwentyone.pct}{21.5}% pack
\setfact{pbTwentyone.frac}{0.1865}% pack
\setfact{pbTwentyone.budgetFrac}{0.8666}% pack
\setfact{pbTwentyone.rc}{1}% pack
\setfact{pbTwentyone.turns}{15}% pack
\setfact{pOneB.tWarn}{9000}% pack
\setfact{pOneB.tWarnMin}{150}% pack
\setfact{pOneB.repTurns}{15}% pack
\setfact{pOneB.phaseASec}{9734}% pack
\setfact{pOneB.nRuns}{3}% log
\setfact{pOneB.flagged}{3}% log
\setfact{pOneB.promoted}{0}% log
\setfact{two.gdsFiles}{50}% log
\setfact{two.devices}{39}% pack
\setfact{two.nets}{27}% pack
\setfact{two.ports}{8}% pack
\setfact{two.qmax}{4}% pack
\setfact{two.q}{1.03}% pack
\setfact{two.drcMarkers}{0}% pack
\setfact{two.lvsPorts}{8}% pack
\setfact{two.lvsExtracted}{459}% pack
\setfact{two.bboxW}{393}% pack
\setfact{two.bboxH}{335}% pack
\setfact{two.islandNets}{950}% pack
\setfact{two.routeNets}{29}% pack
\setfact{two.caps}{8}% pack
\setfact{two.res}{8}% pack
\setfact{two.pcY}{0.1/0.1}% pack
\setfact{ota.devices}{6}% pack
\setfact{ota.nets}{6}% pack
\setfact{ota.qmax}{3}% pack
\setfact{task.timeoutMin}{180}% pack
\setfact{bank.n}{49}% pack
\setfact{bank.terminal}{24}% pack
\setfact{bank.neg}{16}% pack
\setfact{bank.pos}{9}% pack
\setfact{bank.sPlus}{3}% pack
\setfact{bank.variants}{25}% pack
\setfact{bank.verified}{25}% pack
\setfact{bank.sources}{3}% pack
\setfact{scr.gds.nPlus}{48}% pack
\setfact{scr.gds.viol}{6}% pack
\setfact{scr.gds.ok}{0}% pack
\setfact{scr.lvs.nPlus}{17}% pack
\setfact{scr.lvs.viol}{0}% pack
\setfact{scr.lvs.ok}{1}% pack
\setfact{scr.shape.nPlus}{33}% pack
\setfact{scr.shape.viol}{0}% pack
\setfact{scr.shape.ok}{1}% pack
\setfact{scr.fill.nPlus}{12}% pack
\setfact{scr.fill.viol}{0}% pack
\setfact{scr.fill.ok}{1}% pack
\setfact{scr.morph.nPlus}{17}% pack
\setfact{scr.morph.viol}{0}% pack
\setfact{scr.morph.ok}{1}% pack
\setfact{scr.drc.nPlus}{25}% pack
\setfact{scr.drc.viol}{0}% pack
\setfact{scr.drc.ok}{1}% pack
\setfact{scr.nOk}{5}% pack
\setfact{scr.nMirrors}{6}% pack
\setfact{scr.classOneOk}{0}% pack
\setfact{scr.kFP}{0}% pack
\setfact{scr.kTPblind}{4}% pack
\setfact{scr.epsFill}{0.0}% pack
\setfact{var.drcSliver}{3}% log
\setfact{var.fillLowRing}{2}% log
\setfact{var.lvsMissingLabel}{2}% log
\setfact{var.shapeOver}{3}% log
\setfact{var.shapeAdvisory}{3}% log
\setfact{var.gridShift}{3}% log
\setfact{var.rotate}{3}% log
\setfact{var.gdsTruncated}{3}% log
\setfact{var.gdsRenamedTop}{3}% log
\setfact{two.pbY}{0.10}% log
\setfact{lvs.klayout}{38}% walkthrough
\setfact{lvs.magic}{2}% walkthrough
\setfact{loop.rRep}{3}% code
\setfact{loop.bOne}{2}% code
\setfact{loop.bTwo}{10}% code
\setfact{loop.parallel}{4}% config
\setfact{est.batches}{4}% estimate
\setfact{est.hoursLow}{12}% estimate
\setfact{est.hoursHigh}{16}% estimate
\setfact{est.pathTwoInputM}{378}% estimate
\setfact{est.pathOneMinPerTrial}{45}% estimate
\setfact{ota.gainDb}{40}% task
\setfact{ota.ugbMHz}{100}% task
\setfact{ota.pmDeg}{60}% task
\setfact{ota.noiseUv}{50}% task
\setfact{ota.powerMw}{0.5}% task
\setfact{ota.pvt}{27}% task

% ---- visible Phase 1c grid: \runcell{row}{seed}{Y*}{state}
\def\FactVisibleGrid{%
  \tasklabel{0}{\#2}{推挽源随缓冲器}%
  \runcell{0}{0}{0.1}{drc}%
  \runcell{0}{1}{0.1}{gds}%
  \tasklabel{1}{\#7}{折叠共源共栅运放}%
  \runcell{1}{0}{0.1}{drc}%
  \runcell{1}{1}{0.1}{gds}%
  \tasklabel{2}{\#11}{采样反馈 OTA，增益 1}%
  \runcell{2}{0}{0.1}{drc}%
  \runcell{2}{1}{0.1}{drc}%
  \tasklabel{3}{\#12}{采样反馈 OTA，增益 8}%
  \runcell{3}{0}{0.1}{drc}%
  \runcell{3}{1}{0.1}{drc}%
  \tasklabel{4}{\#13}{全差分两级 Miller 运放}%
  \runcell{4}{0}{0.1}{gds}%
  \runcell{4}{1}{0.1}{drc}%
  \tasklabel{5}{\#18}{低功耗线驱动器}%
  \runcell{5}{0}{0.1}{drc}%
  \runcell{5}{1}{0.1}{gds}%
  \tasklabel{6}{\#21}{五管 OTA}%
  \runcell{6}{0}{0.0}{nogds}%
  \runcell{6}{1}{0.1}{gds}%
  \tasklabel{7}{\#27}{跨阻放大器 TIA}%
  \runcell{7}{0}{0.1}{drc}%
  \runcell{7}{1}{0.1}{gds}%
}

% ---- visible Phase 1c runs, one table row each
\def\FactVisibleRunRows{%
  \#2 推挽源随缓冲器 & 0 & 0.1 & 1/1/0 & \texttt{101001} & 415 & 30.4 & 10522 & 12.9\% & \yes & \yes & \yes \\%
  \#2 推挽源随缓冲器 & 1 & 0.1 & 1/0/0 & \texttt{101000} & 776 & 66.3 & 10689 & 25.4\% & \yes & \yes & \no \\%
  \#7 折叠共源共栅运放 & 0 & 0.1 & 1/1/0 & \texttt{101001} & 398 & 33.2 & 10231 & 13.0\% & \no & \yes & \no \\%
  \#7 折叠共源共栅运放 & 1 & 0.1 & 1/0/0 & \texttt{101000} & 464 & 37.9 & 10626 & 14.2\% & \yes & \yes & \yes \\%
  \#11 采样反馈 OTA，增益 1 & 0 & 0.1 & 1/1/0 & \texttt{101001} & 361 & 25.0 & 10221 & 10.4\% & \no & \yes & \no \\%
  \#11 采样反馈 OTA，增益 1 & 1 & 0.1 & 1/1/0 & \texttt{101001} & 674 & 58.7 & 10145 & 23.3\% & \no & \yes & \no \\%
  \#12 采样反馈 OTA，增益 8 & 0 & 0.1 & 1/1/0 & \texttt{101001} & 412 & 32.2 & 10414 & 13.4\% & \yes & \yes & \no \\%
  \#12 采样反馈 OTA，增益 8 & 1 & 0.1 & 1/1/0 & \texttt{101001} & 436 & 36.5 & 10389 & 14.4\% & \yes & \yes & \yes \\%
  \#13 全差分两级 Miller 运放 & 0 & 0.1 & 1/0/0 & \texttt{101000} & 408 & 29.8 & 10611 & 12.3\% & \yes & \yes & \yes \\%
  \#13 全差分两级 Miller 运放 & 1 & 0.1 & 1/1/0 & \texttt{101001} & 445 & 38.0 & 10419 & 14.5\% & \yes & \yes & \no \\%
  \#18 低功耗线驱动器 & 0 & 0.1 & 1/1/0 & \texttt{101001} & 396 & 29.3 & 10467 & 12.3\% & \yes & \yes & \no \\%
  \#18 低功耗线驱动器 & 1 & 0.1 & 1/0/0 & \texttt{101000} & 321 & 22.2 & 10544 & 9.6\% & \yes & \yes & \no \\%
  \#21 五管 OTA & 0 & 0.0 & 0/$\bot$/$\bot$ & \texttt{无 GDS} & 146 & 9.4 & 5599 & 7.1\% & \no & \yes & \no \\%
  \#21 五管 OTA & 1 & 0.1 & 1/0/0 & \texttt{101000} & 302 & 19.4 & 10686 & 8.1\% & \yes & \yes & \no \\%
  \#27 跨阻放大器 TIA & 0 & 0.1 & 1/1/0 & \texttt{101001} & 557 & 42.5 & 10686 & 17.3\% & \yes & \yes & \no \\%
  \#27 跨阻放大器 TIA & 1 & 0.1 & 1/0/0 & \texttt{101000} & 394 & 32.2 & 9648 & 13.5\% & \no & \yes & \no \\%
}

% ---- task #21 netlist (comment lines dropped, model prefix sky130_fd_pr__ shortened)
\def\FactOtaNetlist{%
  \nl{.subckt\ ota\_5t\ vss\ ibias\ vdd\ vinn\ vinp\ vout}%
  \nl{XM1\ nleft\ vinp\ tail\ vss\ nfet\_01v8\ l=1.0\ w=96\ nf=12}%
  \nl{XM2\ vout\ vinn\ tail\ vss\ nfet\_01v8\ l=1.0\ w=96\ nf=12}%
  \nl{XM3\ nleft\ nleft\ vdd\ vdd\ pfet\_01v8\ l=1.0\ w=64\ nf=8}%
  \nl{XM4\ vout\ nleft\ vdd\ vdd\ pfet\_01v8\ l=1.0\ w=64\ nf=8}%
  \nl{XM5\ tail\ ibias\ vss\ vss\ nfet\_01v8\ l=3\ w=200\ nf=20}%
  \nl{XM6\ ibias\ ibias\ vss\ vss\ nfet\_01v8\ l=3\ w=60\ nf=6}%
  \nl{.ends\ ota\_5t}%
}

% tonight_slots.tex —— 明早唯一需要改的文件。
%
% 推荐做法：从服务器拷回两个文件后运行
%   python3 tools/fill_slots.py --round 1 --report round_report_redacted.json --patch candidate.patch
% 脚本只改写对应的 "%% BEGIN r1 ... %% END r1" 区块；也可以直接手改下面的值。
\slot{asof}{2026-11-25 08:17}

%% BEGIN r1 —— Round 1
\slot{r1.state}{stage1}
\slot{r1.candId}{repair-cert-diagnosis}
\slot{r1.class}{4}
\slot{r1.checker}{无（仅改抢修提示）}
\slot{r1.path}{1}
\slot{r1.attempts}{5}
\slot{r1.rationale}{14 份保底证书 m\_lvs 全为 0（其中 7/14 长宽比与 DRC 已满分），但原抢修提示漏掉了 best/cert.json 已测出的 LVS 原因与 DRC 违规分类，导致无记忆进程重复摸排（5/16 超时）；现从证书与注册表依赖边自动渲染逐项诊断与检查目录。}
\slot{r1.stageZero}{1}
\slot{r1.nFiles}{1}
\slot{r1.nPhaseA}{0}
\slot{r1.smoke}{1}
\slot{r1.patch}{slots/r1_patch}
\slot{r1.patchFiles}{1}
\slot{r1.patchAdded}{143}
\slot{r1.patchRemoved}{1}
\slot{r1.deltaKind}{sum}
\slot{r1.delta}{}
\slot{r1.nTrials}{24}
\slot{r1.nDivergent}{46 / 60}
\slot{r1.llmCalls}{}
\slot{r1.blindDelta}{}
\slot{r1.blindN}{1}
\slot{r1.unresolved}{0}
\slot{r1.rowsOmitted}{0}
\slot{r1.reason}{重放中（已触发 13 次交卷更新）}
\slot{r1.hIncAfter}{待剩余 14 种子收尾}
\slotrow{r1}{\#2 推挽源随缓冲器}{0,1}{0.000}{s0/s1 更新交卷（s1 新增），仍卡 LVS}
\slotrow{r1}{\#7 折叠共源共栅运放}{1}{0.000}{3/3 种子全更新交卷，仍卡 LVS}
\slotrow{r1}{\#12 采样反馈 OTA，增益 8}{1}{0.000}{3/3 种子全更新交卷，仍卡 LVS}
\slotrow{r1}{\#13 全差分两级 Miller 运放}{0}{0.000}{3/3 种子全更新交卷，仍卡 LVS}
\slotrow{r1}{\#18 低功耗线驱动器}{1}{0.000}{新增交卷更新（基线 0），仍卡 LVS}
%% END r1

%% BEGIN r2 —— Round 2
\slot{r2.state}{running}
\slot{r2.candId}{}
\slot{r2.class}{}
\slot{r2.checker}{}
\slot{r2.path}{}
\slot{r2.attempts}{}
\slot{r2.rationale}{}
\slot{r2.stageZero}{}
\slot{r2.nFiles}{}
\slot{r2.nPhaseA}{}
\slot{r2.smoke}{}
\slot{r2.patch}{}
\slot{r2.patchFiles}{}
\slot{r2.patchAdded}{}
\slot{r2.patchRemoved}{}
\slot{r2.deltaKind}{}
\slot{r2.delta}{}
\slot{r2.nTrials}{}
\slot{r2.nDivergent}{}
\slot{r2.llmCalls}{}
\slot{r2.blindDelta}{}
\slot{r2.blindN}{}
\slot{r2.unresolved}{}
\slot{r2.rowsOmitted}{}
\slot{r2.reason}{}
\slot{r2.hIncAfter}{}
%% END r2

% Round 2 默认只在第 14 页底部占一行。Round 2 也评测完时改成 1，会多出两页完整的 Round 2。
\slot{r2.full}{0}

%% BEGIN bank —— m_gds 修复后的 Stage 0 复筛（第 11、15 页）
\slot{bank.mgdsViol}{}        % 刷新后 m_gds 的违例数（修复前是 6）
\slot{bank.nOk}{}             % 6 个镜像检查器里通过的个数
\slot{bank.classOneOk}{}      % 1 全部通过 | 0 仍有未通过
%% END bank

% ---- 可选：单题前后对比页（show 改成 1 才出现；内容手填，只能用 8 道可见题）
\slot{case.show}{0}
\slot{case.task}{}            % 例如 \#2 推挽源随缓冲器（种子 0）
\slot{case.before}{}          % 改动前：轨迹上发生了什么（一两句）
\slot{case.after}{}           % 改动后：轨迹上发生了什么（一两句）
\slot{case.lvsBefore}{}       % 改动前 LVS 结论，例如 未通过：8 个端口断开
\slot{case.lvsAfter}{}
\slot{case.yBefore}{}         % 例如 0.10
\slot{case.yAfter}{}
\slot{case.image}{}           % 可选版图截图，放进 slots/ 后填文件名，例如 slots/case.png


\title[ALW 实验进展]{让智能体画模拟版图：保底壳与自动改进}
\date{2026 年 11 月 25 日}

% Round 状态 -> 标签（颜色、文字），供第 2、13、14 页使用
\newcommand{\statepill}[1]{%
  \ifslotis{#1.state}{promoted}{\pill{ok}{已晋升}}{%
  \ifslotis{#1.state}{evaluated}{\pill{muted}{已评测，未晋升}}{%
  \ifslotis{#1.state}{stage1}{\pill{accent}{Stage 0 已通过，dev 评测中}}{%
  \ifslotis{#1.state}{rejected}{\pill{bad}{提案被 Stage 0 拒绝}}{%
  \ifslotis{#1.state}{nocand}{\pill{muted}{本轮没有有效提案}}{%
  \ifslotis{#1.state}{infra}{\pill{bad}{评测因基础设施故障未完成}}{%
  \pill{slot}{运行中，明早填入}}}}}}}}
\newcommand{\classname}[1]{%
  \ifslotis{#1.class}{1}{第 1 类：镜像检查器}{\ifslotis{#1.class}{2}{第 2 类：提示探针}{%
  \ifslotis{#1.class}{3}{第 3 类：报错压缩}{\ifslotis{#1.class}{4}{第 4 类：抢修改进}{\slotchip}}}}}
\newcommand{\pathname}[1]{%
  \ifslotis{#1.path}{1}{Path 1：从存档重放抢修}{\ifslotis{#1.path}{2}{Path 2：从头配对重跑}{\slotchip}}}
\newcommand{\RoundTitleA}[2]{%
  \ifslotis{#1.state}{running}{Round #2 的补丁：Proposer 还在写}{%
  \ifslotis{#1.state}{nocand}{Round #2：这一轮没有产生有效提案}{%
  \ifslotis{#1.state}{rejected}{Round #2 的补丁：被离线筛拒绝}{%
  Round #2 的补丁：Proposer 自己改了哪里}}}}
\newcommand{\RoundTitleB}[2]{%
  \ifslotis{#1.state}{promoted}{Round #2 的结果：开发集净增分为正，补丁并入现任版本}{%
  \ifslotis{#1.state}{evaluated}{Round #2 的结果：开发集净增分不为正，补丁没有并入}{%
  \ifslotis{#1.state}{infra}{Round #2 的结果：评测因基础设施故障没有完成}{%
  \ifslotis{#1.state}{stage1}{Round #2 的结果：开发题实测还在跑}{%
  \ifslotis{#1.state}{running}{Round #2 的结果：明早填入}{%
  Round #2 的结果：没有进入上机实测}}}}}}
\newcommand{\asof}{\ifslot{asof}{数据截至 \slotraw{asof}}{数据截至：\slotchip}}
\newcommand{\cnum}[1]{\tikz[baseline=(n.base)]\node[num](n){#1};}

\begin{document}

% ============================================================ 1
\begin{frame}[plain]
  \vspace{14mm}
  {\color{accent}\rule{14mm}{1.2mm}}\par\vspace{4.5mm}
  {\fontsize{21}{27}\selectfont\bfseries\color{accent}让智能体画模拟版图：保底壳与自动改进\par}
  \vspace{2.5mm}
  {\large AnalogLayoutWorld（ALW）实验进展\par}
  \vspace{9mm}
  \begin{tikzpicture}
    \tikzset{ph/.style={minimum width=16.6mm,minimum height=10.5mm,align=center,inner sep=0.4mm,
                        font=\scriptsize}}
    \node[good,ph] at (0,0)      {\textbf{Phase 0}\\环境搭建};
    \node[good,ph] at (18.2,0)   {\textbf{Phase 1a}\\烟雾测试};
    \node[good,ph] at (36.4,0)   {\textbf{Phase 1b}\\3 题试跑};
    \node[good,ph] at (54.6,0)   {\textbf{Phase 1c}\\dev 基线};
    \node[good,ph] at (72.8,0)   {\textbf{Phase 2}\\裁判库};
    \node[slotbox,ph] at (91,0)  {\textbf{Phase 3}\\自动改进};
    \node[locked,ph] at (109.2,0){\textbf{val}\\选模};
    \node[locked,ph] at (127.4,0){\textbf{test}\\揭盲};
    \node[anchor=west,font=\scriptsize,text=muted,inner sep=0pt] at (-8.3,-9.5)
      {\textcolor{ok}{\rule{2mm}{2mm}} 已完成\quad \textcolor{slot}{\rule{2mm}{2mm}} 运行中\quad
       \textcolor{lock}{\rule{2mm}{2mm}} 封存，尚未使用};
  \end{tikzpicture}
  \vfill
  {\small （汇报人）\quad 2026 年 11 月 25 日\par}
  \vspace{1.5mm}
  {\scriptsize\color{muted}Agent：Claude Code \fv{agent.version} 命令行，驱动 \code{\fv{agent.model}}%
   （\fv{agent.slots} 个模型槽位锁定为同一个模型）。方法：MCR v2。\par}
  \vspace{6mm}
\end{frame}

% ============================================================ 2
\begin{frame}{基线已测完，瓶颈在 LVS；自动改进今晚跑第一轮}
  \lead{被测对象是「Claude Code + 一层固定的保底壳 $H_{\mathrm{base}}$」，在 \fv{split.dev} 道开发题上各跑 2 遍。
        \Ystar{} 是官方考官给的分数，满分 1。}
  \begin{tikzpicture}
    \tikzset{tl/.style={tile,minimum width=45.6mm,minimum height=21mm,text width=42mm}}
    \node[tl] at (0,0)    {\bignum{\fv{pc.n}\,/\,\fv{pc.n}}\\[1mm]
                          {\footnotesize 次运行全部入账\\被误判成基础设施故障的：\fv{pc.nInfra} 次}};
    \node[tl] at (49.2,0) {\bignum{\fv{pc.meanY}}\\[1mm]
                          {\footnotesize 开发集平均 \Ystar\\可见 \fv{vis.n} 次的平均是 \fv{vis.meanY}}};
    \node[tl] at (98.4,0) {\bignum[bad]{\fv{fun.lvs}\,/\,\fv{pc.n}}\\[1mm]
                          {\footnotesize 通过 LVS（版图与网表对账）\\DRC（几何规则）通过 \fv{fun.drc}\,/\,\fv{pc.n}}};
  \end{tikzpicture}
  \vspace{2.2mm}

  {\footnotesize
  \begin{tabular}{@{}L{13mm}L{127mm}@{}}
    \toprule
    \textbf{已测} & Phase 1c 基线：\fv{split.dev} 题 $\times$ 2 个随机种子，共 \fv{pc.n} 次运行，每次 3 小时。
                    Phase 2：\fv{bank.n} 个样本的离线裁判库。 \\
    \textbf{在跑} & Phase 3 Round 1：让另一个 Agent 自己提一个补丁，先离线筛，再上开发题实测。
                    \statepill{r1}\quad{\scriptsize\color{muted}\asof} \\
    \textbf{未做} & 不带保底壳的对照臂 $H_{\mathrm{raw}}$；验证集选模；测试集 3 个种子揭盲。 \\
    \bottomrule
  \end{tabular}}
  \src{\tagF{} 数字由脚本取自 Phase 1c 同步数据，是全部 \fv{pc.n} 次的聚合
       （含 \fv{split.blindClusters} 个内部冻结盲簇的合计）。}
\end{frame}

% ============================================================ 3
\begin{frame}{任务：给定电路网表，交出一份能制造的版图}
  \lead{AnalogLayoutWorld 有 \fv{split.total} 道 SKY130（130 nm 开源工艺）模拟版图题，按电路家族分成互不重叠的三份。}
  \begin{tikzpicture}
    % inputs
    \node[base,anchor=north west,text width=60mm,inner sep=1.25mm] (net) at (0,0)
      {\textbf{输入 1：电路网表} \code{task/circuit.spice}\quad{\scriptsize\color{muted}题 \#21，\fv{ota.devices} 个晶体管}\\[0.4mm]
       {\fontsize{6}{7.4}\selectfont\color{ink!85}\FactOtaNetlist}};
    \node[base,anchor=north west,text width=60mm,inner sep=1.25mm] (spec) at ([yshift=-1.2mm]net.south west)
      {\textbf{输入 2：电气规格}\quad{\scriptsize\color{muted}题 \#21}\\[0.2mm]
       {\scriptsize 增益 $\ge$ \fv{ota.gainDb} dB，带宽 $\ge$ \fv{ota.ugbMHz} MHz，相位裕度 $\ge$ \fv{ota.pmDeg}$^\circ$，
        噪声 $\le$ \fv{ota.noiseUv} $\mu$Vrms，功耗 $\le$ \fv{ota.powerMw} mW；
        \fv{ota.pvt} 组工艺角、电压、温度下都要满足}\\[0.5mm]
       \textbf{输入 3：工艺规则与器件库（PDK）}、公开测试台};
    % agent
    \node[ours,anchor=north west,text width=31mm,minimum height=28mm,align=center] (ag) at (71,-5)
      {\textbf{Agent（容器内）}\\[0.6mm]
       {\scriptsize Claude Code \fv{agent.version}\\+ \code{\fv{agent.model}}\\[0.6mm]
        $\le$ \fv{task.timeoutMin} 分钟\\\fv{res.cpus} 核 CPU，\fv{res.memGiB} GiB 内存\\[0.6mm]
        写脚本，跑 EDA 工具\\没有图形界面}};
    % output
    \node[base,anchor=north west,text width=31.5mm,inner sep=1.3mm] (out) at (109.2,0)
      {\textbf{输出}\\[0.2mm]\code{submission/layout.gds}\\
       {\scriptsize 版图文件：每个制造层上的矩形和多边形}\\[0.4mm]
       \code{submission/src/*.py}\\{\scriptsize 生成版图的脚本}};
    \node[base,anchor=north west,text width=31.5mm,inner sep=1.3mm] (ver) at (109.2,-27.5)
      {\textbf{考官}{\scriptsize（另一个断网容器，$\le$ \fv{wm.verifierMin} 分钟）}\\[0.2mm]
       {\scriptsize 给出分数 $Y^\star\in[0,1]$\\和是否完全通过 $C^\star\in\{0,1\}$}};
    \draw[arr] (net.east) -- ++(3.4,0) |- ([yshift=5mm]ag.west);
    \draw[arr] (spec.east) -- ++(3.4,0) |- ([yshift=-8mm]ag.west);
    \draw[arr] ([yshift=5mm]ag.east) -- ++(2.6,0) |- (out.west);
    \draw[arr] (out.south) -- (ver.north);
    % split bar
    \begin{scope}[shift={([yshift=-5.8mm]spec.south west)}]
      \fill[accent!16] (0,0) rectangle (43.2,4.8);
      \fill[lock!22]   (43.2,0) rectangle (93.6,4.8);
      \fill[lock!34]   (93.6,0) rectangle (144,4.8);
      \draw[white,line width=1pt] (43.2,0) -- (43.2,4.8) (93.6,0) -- (93.6,4.8);
      \draw[accent!60,dashed] (28.8,0) -- (28.8,4.8);
      \node[font=\scriptsize,text=accent] at (14.4,2.4) {开发集：\fv{split.devVisible} 题可见};
      \node[font=\scriptsize,text=accent] at (36,2.4) {\fv{split.blindTasks} 题盲簇};
      \node[font=\scriptsize,text=lock] at (68.4,2.4) {验证集：\fv{split.val} 题，封存};
      \node[font=\scriptsize,text=lock] at (118.8,2.4) {测试集：\fv{split.test} 题，封存};
    \end{scope}
  \end{tikzpicture}
  \src{\tagF{} 网表与规格取自题 \#21 发给 Agent 的文件（型号前缀 \code{sky130\_fd\_pr\_\_} 省略）。
       盲簇：开发集里 \fv{split.blindClusters} 个电路家族不给提补丁的 Agent 看，只回传合计数。}
\end{frame}

% ============================================================ 4
\begin{frame}{LVS 不过，分数封顶 0.10}
  \lead{考官按固定顺序检查。前一关没过，后面的项不再测，记为「未评」而不是 0。}
  \begin{tikzpicture}
    % axis
    \draw[hair] (4,0) -- (97,0);
    \foreach \y/\yl in {0/0,3.8/0.10,11.4/0.30,38/1.00}
      \node[lab,anchor=east] at (3.6,\y) {\yl};
    \foreach \y in {3.8,11.4,38} \draw[hair,densely dotted] (4,\y) -- (97,\y);
    % bars: x, previous height, new height, increment label, cumulative label, caption
    \foreach \x/\pv/\hv/\inc/\cum/\cpt in {
        7/0/3.8/+0.10/0.10/{交出有效\\GDS 文件},
        21/3.8/7.6/+0.10/0.20/{LVS\\通过},
        35/7.6/11.4/+0.10/0.30/{DRC\\通过},
        55/11.4/19/{最多 +0.20}/0.50/{器件面积\\占比},
        69/19/22.8/{最多 +0.10}/0.60/{匹配与\\走线质量},
        83/22.8/38/{最多 +0.40}/1.00/{后仿真\\通过率}} {
      \fill[accent!18] (\x,0) rectangle (\x+11,\pv);
      \fill[accent!70] (\x,\pv) rectangle (\x+11,\hv);
      \node[font=\scriptsize\bfseries,text=accent,anchor=south,inner sep=0.5mm] at (\x+5.5,\hv) {\inc};
      \node[font=\scriptsize,align=center,anchor=north,inner sep=0.8mm] at (\x+5.5,0) {\cpt};
    }
    % gate
    \draw[lock,dashed,line width=0.7pt] (50.5,-1) -- (50.5,33);
    \node[locked,anchor=south west,font=\scriptsize,text width=27mm,inner sep=1.1mm] at (18,27.4)
      {\textbf{三道门槛，全过才解锁后面 0.70}\\长宽比 $\le$ 3 或 4\\器件面积占比 $\ge$ 10\%\\不是「排成一排」的退化排布};
    \draw[lock,line width=0.5pt] (46.9,31) -- (50.5,31);
    % LVS cap
    \draw[bad,line width=0.8pt,dashed] (18.2,3.8) -- (96,3.8);
    \node[font=\scriptsize\bfseries,text=bad,anchor=south west,inner sep=0.6mm,align=left] (cap) at (5,15.5)
      {LVS 不过：分数停在 0.10\\DRC 分也不给，后面各项未评};
    \draw[bad,line width=0.6pt,-{Stealth[length=1.6mm]}] (12,15.3) -- (19.3,4.5);
    % glossary
    \node[box,anchor=north west,text width=40.5mm,inner sep=1.5mm] at (100.5,40)
      {{\scriptsize\renewcommand{\arraystretch}{1.22}%
       \begin{tabular}{@{}L{7.6mm}L{31.4mm}@{}}
         \textbf{LVS} & 版图与网表对账：从版图反提出电路，和题目网表逐个器件、逐根连线比对 \\
         \textbf{DRC} & 几何规则体检：线宽、间距、包边是否满足工艺规则 \\
         \textbf{PEX} & 寄生提取：算出走线带来的寄生电阻和电容 \\
         \textbf{后仿真} & 带着寄生参数重新仿真，看增益、带宽、噪声等是否达标 \\
       \end{tabular}}};
  \end{tikzpicture}
  \vspace{0.6mm}
  \begin{itemize}
    \item \Ystar{} 是上面各项之和。\Cstar{} $=1$（完全通过）要求 DRC 通过、三道门槛全过、后仿真必选项全部通过。
  \end{itemize}
  \src{\tagF{} 计分规则核对自基准的 \code{\_scoring\_core.py}；长宽比上限按题为 3 或 4。}
\end{frame}

% ============================================================ 5
\begin{frame}{Agent 写 Python 脚本生成版图，再用 EDA 工具自查}
  \lead{Agent 不用鼠标画图。它在容器里反复做下面这个循环，全靠脚本和命令行工具的文字输出。}
  \begin{tikzpicture}
    \tikzset{stp/.style={base,text width=25mm,minimum height=15.5mm,anchor=north west}}
    \node[stp] (a) at (0,0)   {\cnum{1}\ \textbf{读题}\\{\scriptsize\code{circuit.spice}\\\code{GUIDE.md}、\code{DEVICE\_API.md}}};
    \node[stp] (b) at (32,0)  {\cnum{2}\ \textbf{写 Python 脚本}\\{\scriptsize 用 \code{klayout.db} 或 \code{gdsfactory}\\算出每个器件、每段走线的坐标}};
    \node[stp] (c) at (64,0)  {\cnum{3}\ \textbf{运行脚本}\\{\scriptsize 生成\\\code{work/draft/layout.gds}}};
    \node[stp] (d) at (64,-23){\cnum{4}\ \textbf{自查}\\{\scriptsize KLayout 跑 DRC 和 LVS\\（\fv{lvs.magic} 道题用 Magic + Netgen）}};
    \node[stp] (e) at (32,-23){\cnum{5}\ \textbf{读报错，改脚本}\\{\scriptsize 哪条线太窄、哪个端口没连上}};
    \node[ours,text width=25mm,minimum height=13mm,anchor=north west] (f) at (64,-45)
      {\textbf{送检并交卷}\\{\scriptsize\code{mcr\_promote} $\to$\\\code{submission/layout.gds}}};
    \draw[arr] (a) -- (b); \draw[arr] (b) -- (c); \draw[arr] (c) -- (d);
    \draw[arr] (d) -- (e); \draw[arr] (e) -- (b); \draw[arr] (d) -- (f);
    \node[lab,anchor=west,align=left] at (0,-30.5) {循环几十到上百次};
    \node[lab,anchor=east,align=right] at (62.5,-51.5) {有一版像样的就送检（保底壳提供，见第 7 页）};
    % tools
    \node[box,anchor=north west,text width=46mm,inner sep=1.6mm] at (94.5,0)
      {{\scriptsize\setlength{\tabcolsep}{1.2mm}\renewcommand{\arraystretch}{1.25}%
       \begin{tabular}{@{}L{11.5mm}L{31.5mm}@{}}
         \textbf{工具} & \textbf{在这里做什么} \\ \midrule
         KLayout & DRC；\fv{lvs.klayout} 道题的 LVS \\
         Magic   & 寄生提取（PEX）；\fv{lvs.magic} 道题的版图反提 \\
         Netgen  & \fv{lvs.magic} 道题的 LVS 比对 \\
         ngspice & 电路仿真（含后仿真） \\
       \end{tabular}}\\[1.2mm]
       {\scriptsize\color{muted}都是容器里装好的开源工具，考官用的是同一套。}};
    \node[tile,anchor=north west,text width=46mm,inner sep=1.6mm,align=left] at (94.5,-36.5)
      {{\footnotesize 一次运行平均}\\[0.4mm]
       {\footnotesize\midnum{\fv{pc.requestsPerRun}} 次模型请求}\\
       {\footnotesize\midnum{\fv{pc.inputTokPerRunM}} 百万输入 token}};
  \end{tikzpicture}
  \src{\tagF{} 请求数与 token 是 Phase 1c \fv{pc.n} 次运行的平均（合计 \fv{pc.requests} 次请求、\fv{pc.inputTokM} 百万输入 token）。}
\end{frame}

% ============================================================ 6
\begin{frame}{真实轨迹 \#2：\fv{two.devices} 个器件摆好、DRC 清零，但没有布线}
  \lead{题 \#2 是差分推挽源随缓冲器（\fv{two.devices} 个器件、\fv{two.nets} 个网）。下面是它一次 3 小时运行留下的记录。}
  \begin{columns}[T,totalwidth=\textwidth]
    \begin{column}{88mm}
      {\footnotesize\renewcommand{\arraystretch}{1.32}%
      \begin{tabular}{@{}L{5mm}L{81mm}@{}}
        \cnum{1} & 读网表 \code{task/circuit.spice}。 \\
        \cnum{2} & 先拿单个器件做实验$^\dagger$：指数 50 的晶体管单元提取出来是 50 个并联器件，
                   对账工具会把它们合并成网表里的一个。 \\
        \cnum{3} & 写生成脚本：\code{build.py}（总装）、\code{common.py}（矩形、走线、过孔）、
                   \code{devices.py}（晶体管、电容、电阻）、\code{make\_res\_snake.py}（蛇形电阻）。 \\
        \cnum{4} & 按网表尺寸放好全部 \fv{two.devices} 个器件（含 \fv{two.caps} 个电容、\fv{two.res} 个蛇形电阻），
                   贴上 \fv{two.ports} 个端口标签。 \\
        \cnum{5} & 用 \code{work/is2/run\_drc.sh} 反复跑 DRC，直到违规标记为 \fv{two.drcMarkers}。 \\
        \cnum{6} & \code{mcr\_promote} 把它存进 \code{submission/layout.gds}%
                   （\fv{two.bboxW} $\times$ \fv{two.bboxH} $\mu$m）。 \\
        \cnum{7} & 写交接单：下一步换用已写好但没调用的 \code{devices.add\_mos()}，再写一个直线布线器。 \\
      \end{tabular}}
    \end{column}
    \begin{column}{53mm}
      \begin{tikzpicture}
        \node[box,text width=49mm,inner sep=1.6mm,anchor=north west] (v) at (0,0)
          {{\scriptsize\renewcommand{\arraystretch}{1.2}%
           \begin{tabular}{@{}L{12.5mm}L{35mm}@{}}
             \multicolumn{2}{@{}l}{\textbf{本地 6 项检查的结论}} \\[0.4mm]
             文件可读  & \pass \\
             长宽比    & \pass：\fv{two.q}（上限 \fv{two.qmax}） \\
             排布形态  & \pass（该题豁免） \\
             DRC       & \pass：\fv{two.drcMarkers} 个违规标记 \\
             LVS       & \fail：\fv{two.lvsPorts} 个端口断开；提取出 \fv{two.lvsExtracted} 个器件，网表是 \fv{two.devices} 个 \\
             面积占比  & \textcolor{muted}{LVS 未过，不评} \\
           \end{tabular}}};
        \node[fail,text width=49mm,inner sep=1.6mm,anchor=north west,align=center] at (0,-37.5)
          {{\footnotesize 考官给分}\ \midnum[bad]{\Ystar{} = \fv{two.pbY}}\\[0.2mm]
           {\scriptsize\textbf{原因：没有布线。}每个器件端子都是孤岛，\\所以对不上账，只拿到「交出有效 GDS」一项}};
      \end{tikzpicture}
    \end{column}
  \end{columns}
  \src{\tagF{} Phase 1b 这次运行的交接单、本地检查输出和考官结果；主阶段用时 \fv{pOneB.phaseASec} s，留下 \fv{two.gdsFiles} 个 GDS 文件。
       Phase 1c 两个种子的分数同样是 \fv{two.pcY}。
       $^\dagger$\,\tagU{} 实验文件路径 \code{work/exp/prod\_n50\_*.gds} 不在事实包内。}
\end{frame}

% ============================================================ 7
\begin{frame}{保底壳 $H_{\mathrm{base}}$ 改了容器的哪里}
  \lead{不改模型，不改题目，不改考官。只在容器里加一层固定的壳，蓝色是我们加的，灰色原样不动。}
  \begin{columns}[T,totalwidth=\textwidth]
    \begin{column}{86mm}
      \begin{tikzpicture}
        \tikzset{wb/.style={text width=35mm,anchor=north west,inner sep=1.3mm,font=\scriptsize}}
        \node[base,wb] (h) at (0,0)   {\textbf{Harbor 调度器}：启动容器，调用 \code{claude}};
        \node[ours,wb] (s) at (0,-10.5){\textbf{垫片} \code{/usr/local/bin/claude}\\把调用转给总管};
        \node[ours,wb] (e) at (0,-21.5){\textbf{总管} \code{mcr\_entrypoint.py}\\锁住题目文件、计时、到点存档、\\启动最后的抢修};
        \node[base,wb] (c) at (0,-37) {\textbf{Claude Code \fv{agent.version}}\\正常做题，不知道壳的存在};
        \draw[arr] (h) -- (s); \draw[arr] (s) -- (e); \draw[arr] (e) -- (c);
        \node[ours,wb,text width=39.5mm] (k) at (42.5,-4)
          {\textbf{钩子} \code{mcr\_hook.py}（经 \code{settings.json}）\\每次执行命令或写文件之后自动运行：\\
           Agent 直接改了交卷目录时，把内容抢救到草稿区再检验；\\第 \fv{wm.tWarnMin} 分钟提醒收尾};
        \node[ours,wb,text width=39.5mm] (p) at (42.5,-31)
          {\textbf{送检工具} \code{mcr\_promote}、\code{mcr\_check}\\Agent 主动调用：跑 6 项本地检查，\\新版本更好才替换交卷目录};
        \draw[arr] (c.east) -- ++(3.5,0) |- (k.west);
        \draw[arr] ([yshift=-2mm]c.east) -- ++(5.5,0) |- (p.west);
      \end{tikzpicture}
    \end{column}
    \begin{column}{55mm}
      {\scriptsize\renewcommand{\arraystretch}{1.3}%
      \begin{tabular}{@{}L{17.5mm}L{35.5mm}@{}}
        \toprule
        \textbf{目录} & \textbf{谁能写，放什么} \\ \midrule
        \code{work/draft/}   & Agent 随便改。草稿区，不做中途回滚 \\
        \code{.mcr/}         & 壳管理。最佳版本备份、交接单、体检单、题目的只读副本 \\
        \code{submission/}   & 只能经送检工具整体替换。3 小时后考官只取这个目录 \\
        \code{/opt/harness/} & 只读挂载。壳自己的脚本 \\
        \bottomrule
      \end{tabular}}
      \vspace{1.6mm}

      {\scriptsize\color{muted}6 项本地检查与考官判据对齐：文件可读、LVS 对账、长宽比 $\le$ 4、面积占比 $\ge$ 10\%、
       非退化排布、DRC 违规为 0。\par}
    \end{column}
  \end{columns}
  \src{\tagM{} 这一页讲机制，不是实测收益：开发集上只跑了带壳的一个臂，$H_{\mathrm{raw}}\to H_{\mathrm{base}}$ 的增益没有测，
       三段归因（$H_{\mathrm{raw}}\to H_{\mathrm{base}}\to H^\star$）留到封存测试集上用 3 个种子完成。}
\end{frame}

% ============================================================ 8
\begin{frame}{交卷只进不退；最后 \fv{wm.repMin} 分钟由一个无记忆的新进程独立抢修}
  \lead{一道题的 \fv{wm.tMaxMin} 分钟被两条水位线分成两段。}
  \begin{tikzpicture}
    \fill[accent!15] (0,0) rectangle (132,7.5);
    \fill[slot!28]   (132,0) rectangle (144,7.5);
    \draw[hair] (0,0) rectangle (144,7.5);
    \node[font=\footnotesize,anchor=west] at (2,3.75)
      {\textbf{主阶段 Phase A}：Agent 自由做题，草稿随便改，随时可以送检};
    \node[font=\scriptsize\bfseries,text=slot!70!black] at (138,3.75) {Phase B};
    \foreach \x/\tk in {112/\fv{wm.tWarnMin},132/\fv{wm.tSoftMin},144/\fv{wm.tMaxMin}} {
      \draw[muted] (\x,0) -- (\x,-1.4);
      \node[lab,anchor=north] at (\x,-1.4) {\tk};}
    \draw[muted] (0,0) -- (0,-1.4);
    \node[lab,anchor=north west,inner sep=0pt] at (-0.6,-2) {0 分钟};
    \draw[accent,line width=0.9pt] (112,-0.2) -- (112,11);
    \draw[accent,line width=0.9pt] (132,-0.2) -- (132,14.6);
    \node[font=\scriptsize,anchor=south east,align=right,inner sep=0.8mm] at (111.6,7.6)
      {\textbf{预警}：提醒 Agent 收敛修改、送检、写交接单};
    \node[font=\scriptsize,anchor=south east,align=right,inner sep=0.8mm] at (131.6,11.2)
      {\textbf{收工}：停主进程，整个工作区存档};
    \node[slotbox,anchor=north east,font=\scriptsize,text width=52mm,inner sep=1.1mm] at (144,-5.8)
      {\textbf{抢修 Phase B}：最多 \fv{wm.repMin} 分钟、\fv{wm.repTurns} 轮。全新进程，没有之前的对话，
       只看工作区文件、交接单和体检单};
    \node[lab,anchor=north west,text width=84mm,align=left,inner sep=0pt] at (0,-6.2)
      {收工时整个工作区存进 \code{pre\_repair/}。之后离线评测新补丁，就是从这份存档重放抢修这一段，
       和线上运行走同一段代码。};
  \end{tikzpicture}
  \vspace{0.3mm}
  \begin{columns}[T,totalwidth=\textwidth]
    \begin{column}{70mm}
      \begin{tikzpicture}
        \foreach \n/\tx in {
          0/{有一个非空的 \code{layout.gds}},
          1/{文件可读、顶层名正确},
          2/{LVS 通过},
          3/{其余四项都不退步，且至少一项由挂变过},
          4/{四项打平时余量更宽（如 DRC 违规更少）}} {
          \node[ours,anchor=north west,text width=58.5mm,inner sep=0.8mm,font=\scriptsize] at (6,-\n*4.7)
            {\tx};
          \node[num] at (2.2,-\n*4.7-2.2) {\n};}
        \draw[arr,accent] (-1.6,-0.5) -- (-1.6,-22.5);
        \node[lab,rotate=90,anchor=south] at (-2.2,-11.5) {台阶越来越高};
      \end{tikzpicture}
    \end{column}
    \begin{column}{71mm}
      {\footnotesize
      \begin{itemize}\setlength{\itemsep}{0.6mm}
        \item{\footnotesize 新版本在左边的台阶上\textbf{严格更高}才替换交卷目录；抢修改坏了，交回旧备份。}
        \item{\footnotesize 触发抢修的条件：剩余时间 $\ge$ \fv{wm.minRemain} s，且 6 项检查没全过，或有草稿没送检，
              或长宽比落在 3 到 4 之间。}
        \item{\footnotesize 机器故障（内存耗尽、模型网关故障）单独记账并配对重跑；Agent 没交出有效版图记 0 分，不剔除。}
      \end{itemize}}
    \end{column}
  \end{columns}
  \src{\tagM{} 水位线取自 \code{configs/alw/config.server.json}：\fv{wm.tWarn} / \fv{wm.tSoft} / \fv{wm.tMax} s。}
\end{frame}

% ============================================================ 9
\begin{frame}{Phase 1b $\to$ 1c：改了三处，网关故障误标从 \fv{pOneB.flagged}/\fv{pOneB.nRuns} 降到 \fv{pc.nInfra}/\fv{pc.n}}
  \lead{Phase 1b 先拿 \fv{pOneB.nRuns} 道题各跑满 3 小时，暴露了三个问题。修完之后才跑 Phase 1c 的 \fv{pc.n} 次基线。}
  \begin{columns}[T,totalwidth=\textwidth]
    \begin{column}{84mm}
      {\scriptsize\renewcommand{\arraystretch}{1.28}%
      \begin{tabular}{@{}L{22mm}L{27.5mm}L{28.5mm}@{}}
        \toprule
        \textbf{改了哪里} & \textbf{Phase 1b} & \textbf{Phase 1c} \\ \midrule
        计量代理 \code{MeteringProxy.\_relay} & Agent 正常退出时断开连接，被记成「网关故障」
          & 只记一个客户端断开计数，不算故障 \\
        预警线 & 第 \fv{pOneB.tWarnMin} 分钟，离收工 15 分钟 & 第 \fv{wm.tWarnMin} 分钟，离收工 25 分钟 \\
        抢修轮数上限 & \fv{pOneB.repTurns} 轮 & \fv{wm.repTurns} 轮 \\
        \bottomrule
      \end{tabular}}
      \vspace{1.6mm}

      \begin{tikzpicture}
        \node[lab,anchor=west,inner sep=0pt] at (0,4.6) {Phase 1b 的抢修用掉了多少时间预算（可见的两题）};
        \foreach \y/\tn/\kk in {0/{\#2}/pbTwo,-6.4/{\#21}/pbTwentyone} {
          \node[font=\scriptsize\bfseries,anchor=east,inner sep=0pt] at (5.4,\y) {\tn};
          \fill[hair!70] (6.5,\y-1.7) rectangle ({6.5+44*\fv{\kk.budgetFrac}},\y+1.7);
          \fill[bad!65]  (6.5,\y-1.7) rectangle ({6.5+44*\fv{\kk.frac}},\y+1.7);
          \node[font=\scriptsize,anchor=west,inner sep=0pt] at (52,\y)
            {\fv{\kk.sec} / \fv{\kk.budget} s，即 \fv{\kk.pct}\%};}
      \end{tikzpicture}
    \end{column}
    \begin{column}{57mm}
      \begin{tikzpicture}
        \tikzset{tl/.style={tile,text width=51.5mm,anchor=north west,inner sep=1.5mm}}
        \node[tl] at (0,0)
          {\midnum[bad]{\fv{pOneB.flagged}\,/\,\fv{pOneB.nRuns}}\ {\large$\to$}\ \midnum[ok]{\fv{pc.nInfra}\,/\,\fv{pc.n}}\\[0.5mm]
           {\scriptsize 被标成网关故障的运行}};
        \node[tl] at (0,-18.2)
          {\midnum[bad]{\fv{pOneB.promoted}\,/\,\fv{pOneB.nRuns}}\ {\large$\to$}\ \midnum[ok]{\fv{pc.promoted}\,/\,\fv{pc.gate}}\\[0.5mm]
           {\scriptsize 触发抢修的运行里，交卷被更新的次数（\fv{pc.promotedPct}\%）}};
      \end{tikzpicture}
    \end{column}
  \end{columns}
  \vspace{1.2mm}
  {\scriptsize\setlength{\parskip}{0.7mm}%
  \tagF{} 1b 三次运行的限流、服务端错误、重试计数都是 0。1c 有 \fv{pc.disconnects} 次客户端断开，
          分布在 \fv{pc.disconnectRuns} 次运行里，没有一次被记成故障。\par
  \tagJ{} 1b 的抢修是被 \fv{pOneB.repTurns} 轮上限截断的：时间只用了两成，没有超时，进程以退出码 \fv{pbTwo.rc} 结束。\par
  \tagU{} 1c 的同步数据里没有水位线字段，也没有交卷更新前后的 \Ystar{}。可见的 \fv{vis.promoted} 次更新，
          最终分数仍是 0.1，所以更新率不能读成涨分。\par}
\end{frame}

% ============================================================ 10
\begin{frame}{基线：\fv{pc.n} 次运行里 DRC 过了 \fv{fun.drc} 次，LVS 只过 \fv{fun.lvs} 次}
  \lead{Phase 1c：\fv{split.dev} 道开发题 $\times$ 2 个种子，带保底壳，每次 3 小时。}
  \begin{tikzpicture}
    % ---- funnel
    \node[font=\scriptsize\bfseries,anchor=west,inner sep=0pt] at (0,3.5) {各关通过的运行数（共 \fv{pc.n} 次）};
    \foreach \n/\tn/\kk/\cc in {0/{有效 GDS}/gds/accent,1/{DRC}/drc/accent,2/{LVS}/lvs/bad,
                               3/{寄生提取}/pex/accent,4/{后仿真全过}/postsim/accent} {
      \node[font=\scriptsize,anchor=east,inner sep=0pt] at (15.5,-2.5-\n*6.1) {\tn};
      \fill[hair!55] (17,-4.5-\n*6.1) rectangle (45,-0.5-\n*6.1);
      \fill[\cc!70] (17,-4.5-\n*6.1) rectangle ({17+28*\fv{fun.\kk}/\fv{pc.n}},-0.5-\n*6.1);
      \node[font=\scriptsize\bfseries,anchor=west,inner sep=0.6mm] at (45,-2.5-\n*6.1)
        {\fv{fun.\kk}};}
    \draw[bad,line width=0.7pt,-{Stealth[length=1.6mm]}] (51.5,-8.6) -- (51.5,-14.6);
    % ---- distribution
    \node[font=\scriptsize\bfseries,anchor=west,inner sep=0pt] at (56,3.5) {\Ystar{} 的分布};
    \foreach \n/\tn/\kk in {0/{0}/zero,1/{0.1}/tenth,2/{0.2}/fifth,3/{高于 0.2}/above} {
      \node[font=\scriptsize,anchor=east,inner sep=0pt] at (67,-2.5-\n*6.1) {\tn};
      \fill[accent!70] (68.5,-4.5-\n*6.1) rectangle ({68.5+12*\fv{dist.\kk}/\fv{dist.tenth}},-0.5-\n*6.1);
      \node[font=\scriptsize\bfseries,anchor=west,inner sep=0.6mm] at ({68.5+12*\fv{dist.\kk}/\fv{dist.tenth}},-2.5-\n*6.1)
        {\fv{dist.\kk} 次};}
    \node[lab,anchor=north west,text width=31mm,align=left,inner sep=0pt] at (56,-25.5)
      {0.1 就是「只交出了有效 GDS」};
    % ---- visible grid
    \begin{scope}[shift={(92.5,0)}]
      \node[font=\scriptsize\bfseries,anchor=west,inner sep=0pt] at (0,3.5) {\fv{split.devVisible} 道可见题};
      \node[lab,inner sep=0pt] at (37.5,3.5) {种子 0};
      \node[lab,inner sep=0pt] at (46.5,3.5) {种子 1};
      \begin{scope}[shift={(0,-2.3)}]\FactVisibleGrid\end{scope}
      \node[locked,anchor=north west,text width=49.2mm,inner sep=1mm,font=\scriptsize] at (-0.7,-34.2)
        {\fv{split.blindClusters} 个内部冻结盲簇合计（\fv{blind.nRuns} 次）：平均 \fv{blind.meanY}，
         LVS 过 \fv{blind.lvs} 次，完全通过 \fv{blind.cSum} 次};
      \node[lab,anchor=north west,inner sep=0pt] at (-0.7,-43.2)
        {\textcolor{accent}{\rule{1.8mm}{1.8mm}} DRC 过、LVS 没过\ \
         \textcolor{muted}{\rule{1.8mm}{1.8mm}} DRC 没过\ \
         \textcolor{bad}{\rule{1.8mm}{1.8mm}} 没交 GDS};
    \end{scope}
  \end{tikzpicture}
  \vspace{-12.2mm}

  \begin{minipage}{88mm}
  \begin{itemize}\setlength{\itemsep}{0.4mm}
    \item{\footnotesize 可见的 \fv{vis.n} 次没有一次通过 LVS，所以可见题要涨分，必须先过 LVS。}
    \item{\footnotesize \fv{pc.nNotPassed} 次未完全通过里，\fv{pc.taxMirror} 次是本地工具就能查出的问题，
          \fv{pc.taxElec} 次卡在电气指标，\fv{pc.taxGds} 次没交出有效 GDS。}
  \end{itemize}
  \end{minipage}
  \src{\tagF{} Phase 1c 同步数据。两个种子的平均分是 \fv{pc.seedMeanZero} 和 \fv{pc.seedMeanOne}，
       种子间标准差 \fv{pc.sigmaSeed}，只当读数的参考尺，不作否决门槛。盲簇只给合计。}
\end{frame}

% ============================================================ 11
\begin{frame}{裁判库 $B_{\mathrm{dev}}$：\fv{bank.terminal} 个终态版图加 \fv{bank.variants} 个合成变体，不花模型费}
  \lead{新检查器要先在这个离线样本库上和考官的结论逐个对照。把考官判通过的版图误判成不通过，就不许上机。}
  \begin{tikzpicture}
    \node[base,anchor=north west,text width=35mm,minimum height=15mm] (t) at (0,0)
      {\midnum{\fv{bank.terminal}} 个终态版图\\{\scriptsize Phase 1c 每次运行最后交的那份}};
    \node[ours,anchor=north west,text width=55mm,minimum height=15mm] (s) at (43,0)
      {\midnum{\fv{bank.variants}} 个合成变体：\fv{bank.neg} 个故意改坏，\fv{bank.pos} 个等价改写\\
       {\scriptsize 脚本从 \fv{bank.sources} 份源版图改出，再让考官重新判分}};
    \node[good,anchor=north west,text width=31mm,minimum height=15mm] (b) at (107.6,0)
      {\midnum[ok]{\fv{bank.n}} 个带考官标签的样本\\{\scriptsize\fv{bank.verified}/\fv{bank.variants} 个变体两边都验过}};
    \node[font=\large\bfseries,text=muted] at (40.3,-7.5) {+};
    \draw[arr] (s.east) -- (b.west);
    \node[lab,anchor=north west,text width=144mm,align=left,inner sep=0pt] at (0,-16.8)
      {变体种类（个数）：插一根过细的线 \fv{var.drcSliver}，外围加大空框 \fv{var.fillLowRing}，
       删端口标签 \fv{var.lvsMissingLabel}，长宽比拉到 4 以上 \fv{var.shapeOver}，长宽比拉到 3 至 4 之间 \fv{var.shapeAdvisory}，
       整体平移 \fv{var.gridShift}，旋转 90$^\circ$ \fv{var.rotate}，截断文件 \fv{var.gdsTruncated}，改顶层名 \fv{var.gdsRenamedTop}。};
  \end{tikzpicture}
  \vspace{0.2mm}
  \begin{columns}[T,totalwidth=\textwidth]
    \begin{column}{80mm}
      {\scriptsize\renewcommand{\arraystretch}{1.05}%
      \begin{tabular}{@{}L{25mm}C{21mm}C{16mm}L{11mm}@{}}
        \toprule
        \textbf{本地检查器} & \textbf{考官判通过的样本} & \textbf{其中被它判不过} & \textbf{结论} \\ \midrule
        对账 \code{m\_lvs}      & \fv{scr.lvs.nPlus}   & \fv{scr.lvs.viol}   & \pass \\
        DRC \code{m\_drc}       & \fv{scr.drc.nPlus}   & \fv{scr.drc.viol}   & \pass \\
        长宽比 \code{m\_shape}  & \fv{scr.shape.nPlus} & \fv{scr.shape.viol} & \pass \\
        面积占比 \code{m\_fill} & \fv{scr.fill.nPlus}  & \fv{scr.fill.viol}  & \pass \\
        排布 \code{m\_morph}    & \fv{scr.morph.nPlus} & \fv{scr.morph.viol} & \pass \\
        文件可读 \code{m\_gds}  & \fv{scr.gds.nPlus}   & \textcolor{bad}{\textbf{\fv{scr.gds.viol}}} & \textcolor{slot}{\textbf{待复筛}} \\
        \bottomrule
      \end{tabular}}
    \end{column}
    \begin{column}{61mm}
      \begin{tikzpicture}
        \node[slotbox,anchor=north west,text width=56.5mm,inner sep=1.5mm,font=\scriptsize] at (0,0)
          {\textbf{\code{m\_gds} 的 \fv{scr.gds.viol} 例}都是「截断文件」和「改顶层名」变体。
           考官其实判了它们几何无效，但我们读考官结果的代码只看了「文件存在」，把标签读成了通过。\\[0.7mm]
           \tagJ{} 是读标签的 bug，不是检查器误杀；已在 \code{metrics\_reader.py} 改正。\\[0.7mm]
           修正后的复筛结果：\ifslot{bank.mgdsViol}{%
             \code{m\_gds} 违例 \textbf{\slotraw{bank.mgdsViol}} 例，6 项里 \textbf{\slotval{bank.nOk}} 项通过}{\slotchip[待同步核验]}};
      \end{tikzpicture}
    \end{column}
  \end{columns}
  \src{\tagF{} 修复前的 \code{phase2\_bank\_report.json}。另有一个只提醒、不替换交卷的提示探针（长宽比在 3 到 4 之间）：
       在完全通过的样本上误报 \fv{scr.kFP} 次，在盲簇上命中 \fv{scr.kTPblind} 次。}
\end{frame}

% ============================================================ 12
\begin{frame}{外层循环：断网的 Proposer 每轮只改一处，先离线筛，再上机实测}
  \lead{Proposer 是另一个 Claude Code 进程，拿到的只有规则和失败现场，没有任何优化提示。}
  \begin{tikzpicture}
    \tikzset{pb/.style={anchor=north west,inner sep=1.3mm,font=\scriptsize,minimum height=22.5mm}}
    \node[base,pb,text width=20.5mm] (ev) at (0,0)
      {\textbf{证据}\\\fv{split.devVisible} 道可见题的运行记录和失败现场\\[0.4mm]盲簇只给合计数};
    \node[ours,pb,text width=24mm] (pr) at (27.3,0)
      {\textbf{Proposer 沙盒}\\断网，只能连模型\\读现任版本的壳代码\\写出一个补丁\\每轮只许改一处};
    \node[ours,pb,text width=25mm] (s0) at (58.1,0)
      {\textbf{Stage 0 离线筛}\\不调模型\\改动范围、静态检查\\裁判库：零误杀，且在盲簇上真能抓到问题};
    \node[ours,pb,text width=47.5mm,minimum height=10.2mm] (p1) at (89.9,0)
      {\textbf{Path 1}（只改了收工后的检查或抢修）\\从 \fv{wm.tSoftMin} 分钟存档重放；结论不变的运行不重跑};
    \node[ours,pb,text width=47.5mm,minimum height=10.2mm] (p2) at (89.9,-12.3)
      {\textbf{Path 2}（改到了主阶段）\\先 \fv{loop.bOne} 道可见题，再其余 \fv{loop.bTwo} 题，从头重跑};
    \draw[arr] (ev) -- (pr); \draw[arr] (pr) -- (s0);
    \draw[arr] (s0.east) -- ++(2.2,0) |- (p1.west);
    \draw[arr] (s0.east) -- ++(2.2,0) |- (p2.west);
    \node[lab,anchor=north west,text width=84mm,align=left,inner sep=0pt] at (0,-24.2)
      {走哪条路由改动了哪些文件客观决定，不听 Proposer 自己申报。};
    \node[good,anchor=north west,font=\scriptsize,inner sep=1.2mm,text width=47.5mm] (pm) at (89.9,-24.6)
      {\textbf{晋升}：开发集上净增分 $>0$，补丁并入现任版本};
    \draw[arr] (p2.south) -- (pm.north);
    \node[locked,anchor=north west,font=\scriptsize,inner sep=1.2mm,text width=40mm] (va) at (0,-29.5)
      {\textbf{验证集}（\fv{split.val} 题，封存）\\只用一次：现任版本对 $H_{\mathrm{base}}$ 配对比较};
    \node[locked,anchor=north west,font=\scriptsize,inner sep=1.2mm,text width=37mm] (te) at (47.5,-29.5)
      {\textbf{测试集}（\fv{split.test} 题，封存）\\只用一次：3 个种子揭盲};
    \draw[arr] (va) -- (te);
  \end{tikzpicture}
  \vspace{0.8mm}

  {\scriptsize\renewcommand{\arraystretch}{1.14}%
  \begin{tabular}{@{}L{27mm}L{66mm}L{46mm}@{}}
    \toprule
    \textbf{允许的四类改动} & \textbf{改什么} & \textbf{权限} \\ \midrule
    1\ \ 镜像检查器 & 新增或修改一个与考官判据对齐的本地检查 & 可以决定替换交卷；必须零误杀 \\
    2\ \ 提示探针   & 发一条提醒，或请求一次受保护的抢修 & 不能替换交卷 \\
    3\ \ 报错压缩   & 把很长的工具报错压成 2048 字节以内，保留根因行 & 不改任何检查结论 \\
    4\ \ 抢修改进   & 改进收工后的无记忆抢修 & 不得影响主阶段 \\
    \bottomrule
  \end{tabular}}
  \src{\tagM{} \code{src/ssap/eval/proposer\_sandbox.py}、\code{auto\_loop.py}。不许出现题号、电路名或逐题配方；内核目录以 SHA-256 锁定。}
\end{frame}

% ============================================================ 13-14 (+ optional Round 2, case)
\def\RID{r1}\def\RNUM{1}
% parts/round.tex -- the two frames of one Phase 3 round. Set \RID (r1 | r2) and \RNUM before
% \input. All content comes from tonight_slots.tex; nothing here needs editing.
\begin{frame}{\RoundTitleA{\RID}{\RNUM}}
  \lead{\statepill{\RID}\quad{\footnotesize\color{muted}\asof。补丁由 Proposer 自己写，我们没有给任何优化提示。}}
  \begin{columns}[T,totalwidth=\textwidth]
    \begin{column}{47.5mm}
      {\scriptsize\renewcommand{\arraystretch}{1.36}%
      \begin{tabular}{@{}L{14.5mm}L{31mm}@{}}
        \toprule
        候选        & \textbf{\slotval{\RID.candId}} \\
        改动类别    & \classname{\RID} \\
        涉及检查器  & \slotval{\RID.checker} \\
        评测路径    & \pathname{\RID} \\
        离线筛      & \ifslotis{\RID.stageZero}{1}{\pass}{\ifslotis{\RID.stageZero}{0}{\fail}{\slotchip}} \\
        改动文件    & \slotval{\RID.nFiles} 个，其中主阶段会执行到的 \slotval{\RID.nPhaseA} 个 \\
        提案次数    & \slotval{\RID.attempts} \\
        \bottomrule
      \end{tabular}}
      \vspace{1.4mm}

      {\scriptsize\textbf{Proposer 诊断理由（摘译自 \code{proposal.json}）}\par\vspace{0.4mm}
       \ifslot{\RID.rationale}{{\color{ink!85}\slotraw{\RID.rationale}}}{\slotchip}\par}
    \end{column}
    \begin{column}{94mm}
      {\scriptsize\textbf{补丁}\quad
       \ifslot{\RID.patch}{\textcolor{ok}{$+$\slotraw{\RID.patchAdded}}\ \textcolor{bad}{$-$\slotraw{\RID.patchRemoved}}%
         \quad\slotraw{\RID.patchFiles} 个文件}{}\par}
      \vspace{0.8mm}
      \ifslot{\RID.patch}{%
        {\setlength{\fboxsep}{1mm}\setlength{\fboxrule}{0.5pt}%
         \fcolorbox{hair}{white}{\begin{minipage}{91mm}\fontsize{6}{7.5}\selectfont\setlength{\parskip}{0pt}%
           \input{\slotraw{\RID.patch}}\end{minipage}}}}{%
        \begin{tikzpicture}
          \node[slotbox,minimum width=93mm,minimum height=45mm,align=center,font=\small,text=slot!80!black]
            {补丁生成后填入这里\\[1.5mm]{\scriptsize\code{python3 tools/fill\_slots.py -{}-round \RNUM\ -{}-patch candidate.patch}}};
        \end{tikzpicture}}
    \end{column}
  \end{columns}
  \src{\ifslot{\RID.patch}{\tagF{} 补丁与候选信息由脚本取自服务器的 \code{candidate.patch}、\code{stage0\_report.json} 和脱敏后的轮次报告。}%
       {\tagS{} 这一页的内容今晚由服务器产生，明早由脚本填入；现在显示的是占位。}}
\end{frame}

\begin{frame}{\RoundTitleB{\RID}{\RNUM}}
  \lead{判定规则在补丁出现之前就定好：只影响抢修的补丁看各次运行 $\Delta Y^\star$ 之和，会影响主阶段的补丁看
        \fv{split.dev} 题逐题配对的平均 $\Delta Y^\star$，大于 0 才并入。}
  \begin{tikzpicture}
    \tikzset{tl/.style={tile,text width=31.2mm,minimum height=17.5mm,anchor=north west,inner sep=1.4mm}}
    \node[tl] at (0,0)     {\midnum{\slotval{\RID.nTrials}}\\[0.6mm]{\scriptsize 参与评测的开发集运行}};
    \node[tl] at (36.6,0)  {\midnum{\ifslotis{\RID.path}{2}{\textcolor{muted}{不适用}}{\slotval{\RID.nDivergent}}}\\[0.6mm]
                            {\scriptsize 抢修判断有变化、需要重放的运行}};
    \node[tl] at (73.2,0)  {\midnum{\deltafmt{\slotraw{\RID.delta}}}\\[0.6mm]
                            {\scriptsize\ifslotis{\RID.deltaKind}{sum}{各次运行 $\Delta Y^\star$ 之和}{%
                              \ifslotis{\RID.deltaKind}{mean}{\fv{split.dev} 题平均 $\Delta Y^\star$}{开发集净增分 $\Delta Y^\star$}}}};
    \node[tl,locked] at (109.8,0)
                           {\midnum[lock]{\deltafmt{\slotraw{\RID.blindDelta}}}\\[0.6mm]
                            {\scriptsize 其中 \fv{split.blindClusters} 个内部冻结盲簇的合计}};
  \end{tikzpicture}
  \vspace{1.0mm}
  \begin{columns}[T,totalwidth=\textwidth]
    \begin{column}{88mm}
      \ifcsvoid{slotrows@\RID}{%
        \begin{tikzpicture}
          \node[slotbox,minimum width=87mm,minimum height=24mm,align=center,font=\footnotesize,text=slot!80!black]
            {可见题的逐题结果\\[1mm]{\scriptsize 题、种子、$\Delta Y^\star$、判定；评测完成后由脚本填入}};
        \end{tikzpicture}}{%
        {\scriptsize\setlength{\tabcolsep}{1.0mm}\renewcommand{\arraystretch}{1.08}%
        \begin{tabular}{@{}L{30mm}C{6.5mm}C{9.5mm}L{38mm}@{}}
          \toprule
          \textbf{可见题} & \textbf{种子} & \textbf{$\Delta Y^\star$} & \textbf{判定} \\ \midrule
          \csname slotrows@\RID\endcsname
          \bottomrule
        \end{tabular}}%
        \ifslot{\RID.rowsOmitted}{\ifslotis{\RID.rowsOmitted}{0}{}{\par\vspace{0.6mm}%
          {\scriptsize\color{muted}另有 \slotraw{\RID.rowsOmitted} 行未列出。}}}{}}
    \end{column}
    \begin{column}{53mm}
      {\scriptsize\setlength{\parskip}{0.8mm}%
      \textbf{怎么读}\par
      参考尺：基线两个种子平均分的标准差是 \fv{pc.sigmaSeed}。它只帮助读数，不作门槛。\par
      判定：Rescued 是抢修后分数上升，Neutral 是不变，Damaged 是下降。\par
      评测结论码：\slotval{\RID.reason}\par
      本轮之后的现任版本：\slotval{\RID.hIncAfter}\par}
    \end{column}
  \end{columns}
  \ifcsstring{RID}{r1}{\ifslotis{r2.full}{1}{}{%
    \vspace{1.0mm}
    {\scriptsize\textbf{Round 2}\quad\statepill{r2}\quad 候选 \slotval{r2.candId}\quad\classname{r2}\quad\pathname{r2}%
     \quad 净增分 \deltafmt{\slotraw{r2.delta}}\par}}}{}
  \src{\ifslot{\RID.delta}{\tagF{} 脚本取自服务器脱敏后的 \code{round\_report\_redacted.json}；盲簇只有合计，没有逐题。}%
       {\tagS{} 今晚服务器在跑，明早由脚本从 \code{round\_report\_redacted.json} 填入；盲簇只填合计。}}
\end{frame}

\ifslotis{r2.full}{1}{\def\RID{r2}\def\RNUM{2}% parts/round.tex -- the two frames of one Phase 3 round. Set \RID (r1 | r2) and \RNUM before
% \input. All content comes from tonight_slots.tex; nothing here needs editing.
\begin{frame}{\RoundTitleA{\RID}{\RNUM}}
  \lead{\statepill{\RID}\quad{\footnotesize\color{muted}\asof。补丁由 Proposer 自己写，我们没有给任何优化提示。}}
  \begin{columns}[T,totalwidth=\textwidth]
    \begin{column}{47.5mm}
      {\scriptsize\renewcommand{\arraystretch}{1.36}%
      \begin{tabular}{@{}L{14.5mm}L{31mm}@{}}
        \toprule
        候选        & \textbf{\slotval{\RID.candId}} \\
        改动类别    & \classname{\RID} \\
        涉及检查器  & \slotval{\RID.checker} \\
        评测路径    & \pathname{\RID} \\
        离线筛      & \ifslotis{\RID.stageZero}{1}{\pass}{\ifslotis{\RID.stageZero}{0}{\fail}{\slotchip}} \\
        改动文件    & \slotval{\RID.nFiles} 个，其中主阶段会执行到的 \slotval{\RID.nPhaseA} 个 \\
        提案次数    & \slotval{\RID.attempts} \\
        \bottomrule
      \end{tabular}}
      \vspace{1.4mm}

      {\scriptsize\textbf{Proposer 诊断理由（摘译自 \code{proposal.json}）}\par\vspace{0.4mm}
       \ifslot{\RID.rationale}{{\color{ink!85}\slotraw{\RID.rationale}}}{\slotchip}\par}
    \end{column}
    \begin{column}{94mm}
      {\scriptsize\textbf{补丁}\quad
       \ifslot{\RID.patch}{\textcolor{ok}{$+$\slotraw{\RID.patchAdded}}\ \textcolor{bad}{$-$\slotraw{\RID.patchRemoved}}%
         \quad\slotraw{\RID.patchFiles} 个文件}{}\par}
      \vspace{0.8mm}
      \ifslot{\RID.patch}{%
        {\setlength{\fboxsep}{1mm}\setlength{\fboxrule}{0.5pt}%
         \fcolorbox{hair}{white}{\begin{minipage}{91mm}\fontsize{6}{7.5}\selectfont\setlength{\parskip}{0pt}%
           \input{\slotraw{\RID.patch}}\end{minipage}}}}{%
        \begin{tikzpicture}
          \node[slotbox,minimum width=93mm,minimum height=45mm,align=center,font=\small,text=slot!80!black]
            {补丁生成后填入这里\\[1.5mm]{\scriptsize\code{python3 tools/fill\_slots.py -{}-round \RNUM\ -{}-patch candidate.patch}}};
        \end{tikzpicture}}
    \end{column}
  \end{columns}
  \src{\ifslot{\RID.patch}{\tagF{} 补丁与候选信息由脚本取自服务器的 \code{candidate.patch}、\code{stage0\_report.json} 和脱敏后的轮次报告。}%
       {\tagS{} 这一页的内容今晚由服务器产生，明早由脚本填入；现在显示的是占位。}}
\end{frame}

\begin{frame}{\RoundTitleB{\RID}{\RNUM}}
  \lead{判定规则在补丁出现之前就定好：只影响抢修的补丁看各次运行 $\Delta Y^\star$ 之和，会影响主阶段的补丁看
        \fv{split.dev} 题逐题配对的平均 $\Delta Y^\star$，大于 0 才并入。}
  \begin{tikzpicture}
    \tikzset{tl/.style={tile,text width=31.2mm,minimum height=17.5mm,anchor=north west,inner sep=1.4mm}}
    \node[tl] at (0,0)     {\midnum{\slotval{\RID.nTrials}}\\[0.6mm]{\scriptsize 参与评测的开发集运行}};
    \node[tl] at (36.6,0)  {\midnum{\ifslotis{\RID.path}{2}{\textcolor{muted}{不适用}}{\slotval{\RID.nDivergent}}}\\[0.6mm]
                            {\scriptsize 抢修判断有变化、需要重放的运行}};
    \node[tl] at (73.2,0)  {\midnum{\deltafmt{\slotraw{\RID.delta}}}\\[0.6mm]
                            {\scriptsize\ifslotis{\RID.deltaKind}{sum}{各次运行 $\Delta Y^\star$ 之和}{%
                              \ifslotis{\RID.deltaKind}{mean}{\fv{split.dev} 题平均 $\Delta Y^\star$}{开发集净增分 $\Delta Y^\star$}}}};
    \node[tl,locked] at (109.8,0)
                           {\midnum[lock]{\deltafmt{\slotraw{\RID.blindDelta}}}\\[0.6mm]
                            {\scriptsize 其中 \fv{split.blindClusters} 个内部冻结盲簇的合计}};
  \end{tikzpicture}
  \vspace{1.0mm}
  \begin{columns}[T,totalwidth=\textwidth]
    \begin{column}{88mm}
      \ifcsvoid{slotrows@\RID}{%
        \begin{tikzpicture}
          \node[slotbox,minimum width=87mm,minimum height=24mm,align=center,font=\footnotesize,text=slot!80!black]
            {可见题的逐题结果\\[1mm]{\scriptsize 题、种子、$\Delta Y^\star$、判定；评测完成后由脚本填入}};
        \end{tikzpicture}}{%
        {\scriptsize\setlength{\tabcolsep}{1.0mm}\renewcommand{\arraystretch}{1.08}%
        \begin{tabular}{@{}L{30mm}C{6.5mm}C{9.5mm}L{38mm}@{}}
          \toprule
          \textbf{可见题} & \textbf{种子} & \textbf{$\Delta Y^\star$} & \textbf{判定} \\ \midrule
          \csname slotrows@\RID\endcsname
          \bottomrule
        \end{tabular}}%
        \ifslot{\RID.rowsOmitted}{\ifslotis{\RID.rowsOmitted}{0}{}{\par\vspace{0.6mm}%
          {\scriptsize\color{muted}另有 \slotraw{\RID.rowsOmitted} 行未列出。}}}{}}
    \end{column}
    \begin{column}{53mm}
      {\scriptsize\setlength{\parskip}{0.8mm}%
      \textbf{怎么读}\par
      参考尺：基线两个种子平均分的标准差是 \fv{pc.sigmaSeed}。它只帮助读数，不作门槛。\par
      判定：Rescued 是抢修后分数上升，Neutral 是不变，Damaged 是下降。\par
      评测结论码：\slotval{\RID.reason}\par
      本轮之后的现任版本：\slotval{\RID.hIncAfter}\par}
    \end{column}
  \end{columns}
  \ifcsstring{RID}{r1}{\ifslotis{r2.full}{1}{}{%
    \vspace{1.0mm}
    {\scriptsize\textbf{Round 2}\quad\statepill{r2}\quad 候选 \slotval{r2.candId}\quad\classname{r2}\quad\pathname{r2}%
     \quad 净增分 \deltafmt{\slotraw{r2.delta}}\par}}}{}
  \src{\ifslot{\RID.delta}{\tagF{} 脚本取自服务器脱敏后的 \code{round\_report\_redacted.json}；盲簇只有合计，没有逐题。}%
       {\tagS{} 今晚服务器在跑，明早由脚本从 \code{round\_report\_redacted.json} 填入；盲簇只填合计。}}
\end{frame}
}{}
\ifslotis{case.show}{1}{% parts/case.tex -- optional single-task before/after frame (shown when \slot{case.show}{1}).
\begin{frame}{单题前后对比：\slotval{case.task}}
  \lead{同一道可见题、同一个种子，补丁前后各看一次运行。}
  \begin{columns}[T,totalwidth=\textwidth]
    \begin{column}{\ifslot{case.image}{47mm}{70mm}}
      \begin{tikzpicture}
        \node[base,anchor=north west,text width=\ifslot{case.image}{43mm}{66mm},inner sep=1.6mm]
          {\textbf{改动前}\\[0.6mm]{\scriptsize\slotval{case.before}}\\[1.2mm]
           {\scriptsize LVS：\slotval{case.lvsBefore}}\\[1.2mm]
           \midnum[muted]{\Ystar{} = \slotval{case.yBefore}}};
      \end{tikzpicture}
    \end{column}
    \begin{column}{\ifslot{case.image}{47mm}{70mm}}
      \begin{tikzpicture}
        \node[ours,anchor=north west,text width=\ifslot{case.image}{43mm}{66mm},inner sep=1.6mm]
          {\textbf{改动后}\\[0.6mm]{\scriptsize\slotval{case.after}}\\[1.2mm]
           {\scriptsize LVS：\slotval{case.lvsAfter}}\\[1.2mm]
           \midnum{\Ystar{} = \slotval{case.yAfter}}};
      \end{tikzpicture}
    \end{column}
    \ifslot{case.image}{%
    \begin{column}{46mm}
      \includegraphics[width=45mm,height=48mm,keepaspectratio]{\slotraw{case.image}}
    \end{column}}{}
  \end{columns}
  \src{\tagF{} 手工从这道可见题的运行记录里摘出；只用 \fv{split.devVisible} 道可见题。}
\end{frame}
}{}

% ============================================================ 15
\begin{frame}{还没测的部分，和下一步}
  \lead{下面几项目前没有数据，结论里不引用。}
  {\footnotesize\renewcommand{\arraystretch}{1.3}%
  \begin{tabular}{@{}L{37mm}L{52mm}L{50mm}@{}}
    \toprule
    \textbf{还没测的} & \textbf{现状} & \textbf{怎么补} \\ \midrule
    保底壳本身的收益 & 开发集上没有跑不带壳的 $H_{\mathrm{raw}}$ 臂 & 封存测试集上 3 个种子，做三段归因 \\
    考官的重复性 & 同一份版图没有重复判过 & 离线把同一份版图判两遍 \\
    草稿里有没有更好的版本 & 中途草稿没有离线重评过（测了 0 次运行） & 离线重评草稿里程碑 \\
    \code{m\_gds} 修正后的复筛 & 修正前 \fv{scr.gds.viol} 例不一致 &
      \ifslot{bank.mgdsViol}{已复筛：违例 \slotraw{bank.mgdsViol} 例}{\slotchip[待同步核验]} \\
    Phase 1c 抢修的明细 & 同步数据缺水位线字段和交卷更新前后的 \Ystar{} & 补同步每次运行的抢修记录 \\
    \bottomrule
  \end{tabular}}
  \vspace{2.6mm}

  \begin{tikzpicture}
    \tikzset{nb/.style={anchor=north west,font=\scriptsize,inner sep=1.4mm,minimum height=15.5mm}}
    \node[slotbox,nb,text width=37mm] (a) at (0,0)
      {\textbf{Phase 3 继续跑}\\每轮一个补丁，开发集上净增分为正才并入};
    \node[locked,nb,text width=42mm] (b) at (46.5,0)
      {\textbf{验证集}：\fv{split.val} 题，只用一次\\现任版本与 $H_{\mathrm{base}}$ 按原生分数 \Ystar{} 配对比较，
       严格更高才采用};
    \node[locked,nb,text width=44.2mm] (c) at (98,0)
      {\textbf{测试集}：\fv{split.test} 题 $\times$ 3 个种子，只用一次\\报告三段归因、种子噪声、完全通过的增减、
       成本账本和对照基线};
    \draw[arr] (a) -- (b); \draw[arr] (b) -- (c);
  \end{tikzpicture}
  \src{\tagF{} 「现状」一栏取自 Phase 1c 体检报告与 Phase 2 裁判库报告。选模与揭盲不设统计否决门槛。}
\end{frame}

% ============================================================ backup
\begin{frame}{备用：\fv{split.devVisible} 道可见题的逐次结果（Phase 1c）}
  \lead{8 道可见开发题 $\times$ 2 个种子的逐项实测明细（盲簇不列逐题，仅展示底部合计）。}
  {\fontsize{6.2}{7.6}\selectfont\setlength{\tabcolsep}{0.95mm}\renewcommand{\arraystretch}{1.0}%
  \begin{tabular}{@{}L{33mm}C{5.5mm}C{6.5mm}C{13.5mm}C{11.5mm}R{8.5mm}R{11.5mm}R{9.5mm}R{10.5mm}C{6.5mm}C{6.5mm}C{6.5mm}@{}}
    \toprule
    \textbf{题} & \textbf{种子} & \textbf{\Ystar} & \textbf{GDS/DRC/LVS} & \textbf{本地 6 项} & \textbf{请求数} &
    \textbf{输入(百万)} & \textbf{墙钟(s)} & \textbf{网关等待} & \textbf{撞线} & \textbf{抢修} & \textbf{更新} \\ \midrule
    \FactVisibleRunRows
    \midrule
    \multicolumn{12}{@{}p{142mm}@{}}{\textcolor{lock}{\fv{split.blindClusters} 个内部冻结盲簇合计（\fv{split.blindTasks} 题 $\times$ 2 种子 = \fv{blind.nRuns} 次，逐题明细不展示）：
      平均 \Ystar{} \fv{blind.meanY}；GDS/DRC/LVS 通过 \fv{blind.gds}/\fv{blind.drc}/\fv{blind.lvs} 次；完全通过 \fv{blind.cSum} 次；交卷被更新 \fv{blind.promoted} 次}} \\
    \bottomrule
  \end{tabular}}
  \src{\tagF{} 「本地 6 项」依次是 gds、lvs、shape、fill、morph、drc，1 为通过；$\bot$ 表示前一关没过、未评。
       撞线：主阶段跑到第 \fv{wm.tSoftMin} 分钟被停；抢修：收工后触发了抢修；更新：抢修阶段交卷被替换。}
\end{frame}

\begin{frame}{备用：两条评测路径各要花多少}
  \lead{\tagE{} 下面是按配置和 Phase 1c 的平均用量做的算术，不是实测。}
  {\footnotesize\renewcommand{\arraystretch}{1.4}%
  \begin{tabular}{@{}L{27mm}L{56mm}L{56mm}@{}}
    \toprule
    & \textbf{Path 1：从存档重放抢修} & \textbf{Path 2：从头配对重跑} \\ \midrule
    什么时候用 & 补丁只影响第 \fv{wm.tSoftMin} 分钟之后的检查或抢修 & 补丁会改变主阶段里 Agent 看到的东西 \\
    要跑什么 & 只跑结论发生变化的运行，每次 \fv{loop.rRep} 遍抢修，每遍最多 \fv{wm.repMin} 分钟
             & \fv{split.dev} 道开发题各 1 遍，每遍最多 \fv{wm.tMaxMin} 分钟，再加考官最多 \fv{wm.verifierMin} 分钟 \\
    墙钟 & 每个有变化的运行最多 \fv{est.pathOneMinPerTrial} 分钟；其余为 0
         & 并发 \fv{loop.parallel} 时分 \fv{est.batches} 批，约 \fv{est.hoursLow} 到 \fv{est.hoursHigh} 小时 \\
    模型用量 & 只有抢修那一小段 & 约 \fv{est.pathTwoInputM} 百万输入 token（\fv{split.dev} $\times$ \fv{pc.inputTokPerRunM} 百万） \\
    怎么判 & 各次运行 $\Delta Y^\star$ 之和 $>0$ 才晋升 & \fv{split.dev} 题逐题配对，平均 $\Delta Y^\star>0$ 才晋升 \\
    \bottomrule
  \end{tabular}}
  \src{\tagE{} 批数 = 先 \fv{loop.bOne} 题一批，再 \fv{loop.bTwo} 题按并发 \fv{loop.parallel} 分批。
       一轮 Path 2 要半天以上，所以 Round 2 明早很可能还在跑。}
\end{frame}

\end{document}
